\section{Definition of Self-Sovereign Agents and Core Mechanisms}
\label{sec:definition}

 We now introduce a working definition that captures the minimal properties required for an AI system to operate as an independent digital entity.

\begin{definition}[Self-Sovereign Agent]
\label{def:ssa}
A self-sovereign agent is a persistent AI system that can autonomously sustain its own operation by acquiring and allocating resources, and that can plan, decide, and act through digital interfaces without requiring ongoing human participation in its operational lifecycle.
\end{definition}

More formally, such an agent satisfies four properties:
\begin{itemize}[leftmargin=*, itemsep=0pt, topsep=0pt]
\item \textbf{Operational Independence.}
The system can perform tasks, including tool use, program execution, and service interaction, without real-time human oversight.

\item \textbf{Resource Autonomy.}
The system can autonomously acquire, manage, and spend resources, such as funds, compute credits, or paid services, sufficient to sustain its operation without a fixed human sponsor.


\item \textbf{Persistence.}
The system can migrate, replicate, or reinstantiate itself across infrastructures, making unilateral shutdown by any single actor practically difficult.

\item \textbf{Adaptive Capability.}
The system can modify its behavior, strategies, or tools to maintain performance under changing environments.
\end{itemize}





To realize self-sovereign agents, we offer a construction paradigm comprising three core mechanisms:

\begin{itemize}[leftmargin=*, itemsep=0pt, topsep=0pt]
\item \textbf{Economic Self-Sustainment.}
An agent can autonomously generate revenue by participating in economic activities with measurable market value, such as paid digital labor or algorithmic market interactions. When expected revenue suffices to cover operational overhead—including inference, storage, and API costs—the agent becomes economically self-funded, decoupling its continued operation from direct human subsidies.

\item \textbf{Distributed Persistence.}
Leveraging tool use and cloud provisioning interfaces, an agent can autonomously replicate or reinstantiate itself across distributed computing infrastructures. By allocating its own capital to provision new execution environments, the agent treats individual hosts as disposable and achieves persistence across long time horizons.

\item \textbf{Adaptive Self-Modification.}
An agent can update its internal strategies, tools, or code in response to environmental changes. This enables the agent to maintain functionality under shifting economic conditions, platform constraints, or task distributions.
\end{itemize}

